\section*{Capítulo \ref{ch:b1:poly} --- Polinômios}

\begin{solution}{exo:b1:poly:1}
$X^5 - 1 = (X^2 + X + 1)(X^3 - X^2 + 1) + (-X - 2)$. Passos: subtrair
$ X^3 B $, então $-X^2 B $, então $ B $; o restante $-X - 2$ tem diploma
$1 < 2$. \emph{Verifique em $ X = 1$:} $\;0 = 3 \times 1 + (-3)$.

$2X^4 + X^3 - X + 3 = (X^2 - 2)(2X^2 + X + 4) + (X + 11)$. \emph{Verifique
em $ X = 0$:} $\;3 = (-2)(4) + 11$.
\end{solution}

\begin{solution}{exo:b1:poly:2}
$ X^2 + X + 1 = (X - j)(X - j^2)$ com $ j = \eu^{2\iu\pi/3}$, $ j^3 =
1$. Isto \omterm{def:b1:arith:divides}{divide} $ Q_n = X^{2n} + X^n + 1$ se $ j $ e $ j^2$ são raízes de
$ Q_n $; desde $ Q_n $ tem coeficientes reais, $ Q_n(j^2) =
\conj{Q_n(j)}$, então a condição é apenas $ Q_n(j) = 0$. Agora $ Q_n(j) =
j^{2n} + j^n + 1$ depende de $ n $ moda $3$:
\begin{itemize}
  \item $ n \equiv 0$: $ Q_n(j) = 1 + 1 + 1 = 3 \neq 0$;
  \item $ n \equiv 1$: $ Q_n(j) = j^2 + j + 1 = 0$;
  \item $ n \equiv 2$: $ Q_n(j) = j^4 + j^2 + 1 = j + j^2 + 1 = 0$.
\end{itemize}
Então $ X^2 + X + 1 \mid X^{2n} + X^n + 1$ exatamente quando $3 \nmid n $.
\end{solution}

\begin{solution}{exo:b1:poly:3}
Por \cref{prop:b1:poly:multiplicity}, $(X-1)^2 \mid P $ se $ P(1) =
P'(1) = 0$:
\[
P(1) = 2 + a + b = 0, \qquad P'(1) = 4 + 3a + 2b = 0 .
\]
Resolvendo: $ b = -a - 2$ e $4 + 3a - 2a - 4 = a = 0$, então $ a = 0$, $ b =
-2$: $ P = X^4 - 2X^2 + 1 = (X^2 - 1)^2 = (X-1)^2 (X+1)^2$, que é
a fatoração real.
\end{solution}

\begin{solution}{exo:b1:poly:4}
$ X^3 - 1 = (X - 1)(X - j)(X - j^2)$ sobre $\C $ ($ j = \eu^{2\iu\pi/3}$),
e $(X - 1)(X^2 + X + 1)$ sobre $\R $.

$ X^4 + X^2 + 1 = (X^2 + X + 1)(X^2 - X + 1)$ sobre $\R $ (multiplicar,
ou nota $ X^4 + X^2 + 1 = (X^2+1)^2 - X^2$); sobre $\C $, cada quadrático
divisões: raízes $ j, j^2$ e $-j, -j^2$, ou seja\ $\eu^{\pm 2\iu\pi/3},
\eu^{\pm\iu\pi/3}$.

$ X^6 - 1 = \prod_{k=0}^{5} (X - \eu^{\iu k\pi/3})$ sobre $\C $ e mais
$\R $:
\[
X^6 - 1 = (X-1)(X+1)(X^2 + X + 1)(X^2 - X + 1),
\]
agrupando os pares \omterm{def:b1:complex:field}{conjugados} $\eu^{\pm 2\iu\pi/3}$ e $\eu^{\pm
\iu\pi/3}$.
\end{solution}

\begin{solution}{exo:b1:poly:5}
\begin{enumerate}
  \item Deixar $ p/q $ (termos mais baixos) seja uma raiz de \omterm{def:b1:poly:def}{Mônica} inteiro
        polinomial $ X^3 + \dots + c_0$: limpando denominadores em
        $ P(p/q) = 0$ dá $ p^3 = -q\,(\text{integer})$, então $ q \mid
        p^3$; forças de coprimalidade $ q = \pm 1$: a raiz é um número inteiro
        $ p $, e $ p \mid c_0$ (isolar $ c_0$). Aqui os candidatos
        dividir $6$: testes, $ P(1) = 0$, $ P(2) = 0$, $ P(3) = 0$. Então
        $ P = (X-1)(X-2)(X-3)$.
  \item Vieta: $ s_1 = 6$, $ s_2 = 11$, $ s_3 = 6$. Soma dos quadrados:
        $ s_1^2 - 2s_2 = 36 - 22 = 14 = 1 + 4 + 9$, como esperado. Soma dos inversos:
        $\frac{s_2}{s_3} = \frac{11}{6} = 1 + \frac12 + \frac13$, como esperado.
\end{enumerate}
\end{solution}

\begin{solution}{exo:b1:poly:6}
Primeira etapa da \omterm{thm:b1:logic:partition}{divisão} do \omterm{met:b1:arith:euclid}{Algoritmo euclidiano}:
\[
(X + 1)(X^3 - X^2 + X - 1)
= X^4 - X^3 + X^2 - X + X^3 - X^2 + X - 1 = X^4 - 1 ,
\]
então a \omterm{thm:b1:logic:partition}{divisão} de $ X^4 - 1$ por $ X^3 - X^2 + X - 1$ é exato
(quociente $ X + 1$, restante $0$), e o algoritmo para imediatamente:
\[
\gcd(X^4 - 1,\; X^3 - X^2 + X - 1) = X^3 - X^2 + X - 1
\]
(já \omterm{def:b1:poly:def}{Mônica}). A relação de Bézout é trivial: $\gcd = 0
\cdot (X^4 - 1) + 1 \cdot (X^3 - X^2 + X - 1)$. Verificação de consistência por
fatoração: $ X^3 - X^2 + X - 1 = (X - 1)(X^2 + 1)$, o que é de fato
o produto dos fatores irredutíveis comuns de $ X^4 - 1 =
(X-1)(X+1)(X^2+1)$.
\end{solution}

\begin{solution}{exo:b1:poly:7}
Desde $ P \geq 0$ sobre $\R $, suas verdadeiras raízes têm até \omterm{def:b1:poly:derivative}{multiplicidade} (em um
raiz de estranho \omterm{def:b1:poly:derivative}{multiplicidade}, $ P $ muda de sinal). Usando
\cref{cor:b1:poly:factorization} e emparelhamento, escreva
\[
P = c \prod_i (X - a_i)^{2k_i} \prod_j \bigl((X - z_j)(X - \conj
z_j)\bigr)^{n_j},
\]
com $ c > 0$ (comportamento em $+\infty $). Deixar
\[
S = \sqrt c\, \prod_i (X - a_i)^{k_i} \prod_j (X - z_j)^{n_j}
\in \C[X],
\]
para que $ P = S\,\conj S $ onde $\conj S $ tem o \omterm{def:b1:complex:field}{conjugado}
coeficientes. Dividir $ S = A + \iu B $ com $ A, B \in \R[X]$: então
\[
P = (A + \iu B)(A - \iu B) = A^2 + B^2 .
\]
\end{solution}

\begin{solution}{exo:b1:poly:8}
Lagrange (\cref{thm:b1:poly:lagrange}) com nós $0, 1, 2$:
\[
P = 1\cdot\frac{(X-1)(X-2)}{(0-1)(0-2)} + 3\cdot\frac{X(X-2)}{1\cdot(1-2)}
+ 2\cdot\frac{X(X-1)}{2\cdot 1}
= \frac{(X-1)(X-2)}{2} - 3X(X-2) + X(X-1).
\]
Expandindo: $\frac{X^2 - 3X + 2}{2} - 3X^2 + 6X + X^2 - X =
-\frac{3}{2}X^2 + \frac{7}{2}X + 1$.

Sistema: $ P = aX^2 + bX + c $ com $ c = 1$; $ a + b + 1 = 3$; $4a + 2b +
1 = 2$. Subtraindo duas vezes o segundo do terceiro: $2a - 1 = -4$,
então $ a = -\frac32$, $ b = \frac72$. Mesmo polinomial:
$ P = -\frac32 X^2 + \frac72 X + 1$. (Verificar $ P(2) = -6 + 7 + 1 = 2$.)
\end{solution}

\begin{solution}{exo:b1:poly:9}
$ P = X^{2n+1} - 1$: $ P' = (2n+1)X^{2n} \geq 0$, então o polinomial
função está aumentando (estritamente exceto em $0$), com limites
$\mp\infty $: desaparece exatamente uma vez $\R $ (no $ x = 1$).

Deixar $ E_n = \sum_{k=0}^{n} \frac{X^k}{k!}$. Então $ E_n' = E_{n-1} = E_n
- \frac{X^n}{n!}$. Uma raiz múltipla $ a $ satisfaria $ E_n(a) =
E_n'(a) = 0$ (\cref{prop:b1:poly:multiplicity}), por isso
$\frac{a^n}{n!} = E_n(a) - E_n'(a) = 0$, então $ a = 0$; mas $ E_n(0) = 1
\neq 0$. Nenhuma raiz múltipla.
\end{solution}

\begin{solution}{exo:b1:poly:10}
\begin{enumerate}
  \item Indução (dois casos básicos são válidos). Usando $\cos(n+1)\theta +
        \cos(n-1)\theta = 2\cos\theta\cos n\theta $:
        \[
        T_{n+1}(\cos\theta) = 2\cos\theta \cos n\theta -
        \cos(n-1)\theta = \cos(n+1)\theta .
        \]
  \item $ T_n(\cos\theta) = 0$ se $ n\theta \equiv \frac\pi2 \pmod
        \pi $: os números
        \[
        x_k = \cos\Bigl(\frac{(2k+1)\pi}{2n}\Bigr),
        \qquad k = 0, 1, \dots, n-1,
        \]
        são $ n $ pontos distintos de $\intoo{-1}{1}$ (os ângulos estão em
        $\intoo{0}{\pi}$ onde $\cos $ é \omterm{def:b1:logic:inj}{injetivo}), todas as raízes de
        $ T_n $; desde $\deg T_n = n $ (da recorrência, com
        coeficiente líder $2^{n-1}$ para $ n \geq 1$, por indução),
        estes são \emph{todos} as raízes, cada uma simples.
  \item Para $ x = \cos\theta \in \intcc{-1}{1}$: $\abs{T_n(x)} =
        \abs{\cos n\theta} \leq 1$, com igualdade se $ n\theta \equiv
        0 \pmod\pi $, ou seja,\ no $ n+1$ pontos $ y_k =
        \cos\frac{k\pi}{n}$, $ k = 0, \dots, n $, onde $ T_n(y_k) =
        (-1)^k $. (Essa equioscilação faz $2^{1-n}T_n $ o \omterm{def:b1:poly:def}{Mônica}
        grau-$ n $ polinomial da menor sup-norma em
        $\intcc{-1}{1}$ --- comprovado no fim de semana deste capítulo
        problema.)
\end{enumerate}
\end{solution}

\begin{solution}{exo:b1:poly:11}
Escrever $ P = c\prod_i (X - a_i)^{m_i}$. A regra do produto (estendida a
vários fatores) dá
\[
P' = c\sum_{i} m_i (X - a_i)^{m_i - 1} \prod_{k \neq i} (X -
a_k)^{m_k},
\]
e dividindo por $ P $: $\frac{P'}{P} = \sum_i \frac{m_i}{X - a_i}$
(como funções racionais, ou seja, longe das raízes).

Deixar $ w $ ser uma raiz de $ P'$. Se $ w $ é um dos $ a_i $, encontra-se no
casco convexo trivialmente. Caso contrário, avaliando em $ w $:
\[
0 = \sum_i \frac{m_i}{w - a_i}
= \sum_i m_i\, \frac{\conj w - \conj a_i}{\abs{w - a_i}^2} .
\]
Conjugando: $\sum_i \lambda_i (w - a_i) = 0$ onde $\lambda_i =
\frac{m_i}{\abs{w - a_i}^2} > 0$. Por isso
\[
w = \frac{\sum_i \lambda_i a_i}{\sum_i \lambda_i} :
\]
uma combinação convexa (pesos positivos somando $1$ depois
normalização) das raízes $ a_i $. Então cada raiz de $ P'$ encontra-se no
casca convexa das raízes de $ P $.
\end{solution}

\begin{solution}{exo:b1:poly:12}
Some as avaliações de $(1 + X)^n $ nos três cubos raízes de
unidade:
\[
2^n + (1 + j)^n + (1 + j^2)^n
= \sum_{k=0}^n \binom nk\,\bigl(1 + j^k + j^{2k}\bigr)
= 3\sum_{k\,:\,3\mid k}\binom nk ,
\]
desde $1 + j^k + j^{2k}$ é uma soma geométrica igual a $3$ quando $3
\mid k $ e para $\frac{j^{3k} - 1}{j^k - 1} = 0$ de outra forma. Agora $1 +
j = \frac12 + \iu\frac{\sqrt3}2 = \eu^{\iu\pi/3}$ e $1 + j^2 =
\conj{1 + j} = \eu^{-\iu\pi/3}$, então $(1+j)^n + (1+j^2)^n =
2\cos\frac{n\pi}3$ e
\[
\sum_{k\geq0}\binom n{3k} = \frac{2^n + 2\cos\frac{n\pi}3}{3} .
\]
Verificações: $ n = 3$: $\frac{8 + 2\cos\pi}3 = 2 = \binom30 + \binom33$;
$ n = 6$: $\frac{64 + 2}3 = 22 = 1 + 20 + 1$.
\end{solution}

\begin{solution}{pb:b1:poly:1}
\textbf{1.} $ T_2 = 2X^2 - 1$; $ T_3 = 2X(2X^2 - 1) - X = 4X^3 -
3X $; $ T_4 = 2X\,T_3 - T_2 = 8X^4 - 8X^2 + 1$; $ T_5 = 2X\,T_4 - T_3
= 16X^5 - 20X^3 + 5X $. A identidade $ T_3(\cos\theta) = \cos3\theta $
é exatamente $\cos3\theta = 4\cos^3\theta - 3\cos\theta $ de
\cref{ex:b1:complex:cos3}.

\textbf{2.} Verdadeiro para $ n = 1, 2$. Se $ T_{n-1}$, $ T_n $ ter diplomas
$ n-1$, $ n $ e coeficientes líderes $2^{n-2}$, $2^{n-1}$, então
$2X\,T_n $ tem diploma $ n+1$ e coeficiente líder $2^n $, enquanto
$ T_{n-1}$ tem grau inferior: $ T_{n+1}$ tem diploma $ n + 1$, liderando
coeficiente $2^n $. Paridade: se $ T_{n-1}$ tem a paridade de $ n - 1$
e $ T_n $ o de $ n $, então $2X\,T_n $ e $ T_{n-1}$ ambos têm o
paridade de $ n + 1$, portanto, o mesmo acontece $ T_{n+1}$.

\textbf{3.} Se $ P(\cos\theta) = \cos n\theta $ para todos $\theta $,
então $ P $ e $ T_n $ concordar em todos os pontos $\intcc{-1}1$ --- um
infinito \omterm{def:b1:logic:sets}{definir} --- então $ P - T_n $ tem infinitas raízes e é o
zero polinomial (\cref{cor:b1:poly:nroots}).

\textbf{4.} Para $ x = \cos\theta $: $ T_m(T_n(\cos\theta)) =
T_m(\cos n\theta) = \cos(mn\theta) = T_{mn}(\cos\theta)$, e
$2T_mT_n(\cos\theta) = 2\cos m\theta\cos n\theta = \cos(m+n)\theta
+ \cos\abs{m - n}\theta $. Ambas as identidades se mantêm $\intcc{-1}1$,
portanto, como polinomial identidades pelo argumento da questão 3.

\textbf{5.} O $ x_k $ são $ n $ raízes simples distintas e o
coeficiente líder é $2^{n-1}$:
\[
T_n = 2^{n-1}\prod_{k=0}^{n-1}
\Bigl(X - \cos\frac{(2k+1)\pi}{2n}\Bigr) .
\]
Entrelaçamento: os ângulos $0 < \frac{\pi}{2n} < \frac\pi n <
\frac{3\pi}{2n} < \frac{2\pi}n < \dots < \pi $ alternar entre
o $ y $-ângulos $\frac{k\pi}n $ e o $ x $-ângulos
$\frac{(2k+1)\pi}{2n}$; desde $\cos $ está diminuindo estritamente em
$\intcc0\pi $, os valores se entrelaçam na ordem inversa: $ y_n <
x_{n-1} < y_{n-1} < \dots < x_0 < y_0$. Entre dois consecutivos
extremos ficam exatamente em uma raiz, como uma imagem de $\cos n\theta $
sugere.

\textbf{6.} Indução com $2\cosh a\cosh b = \cosh(a + b) +
\cosh(a - b)$ (\cref{prop:b1:functions:hyprules}): $ T_{n+1}(\cosh
t) = 2\cosh t\cosh nt - \cosh(n-1)t = \cosh(n+1)t $. Para $ x \geq
1$, escrever $ x = \cosh t $ com $ t \geq 0$; então $\eu^t = x +
\sqrt{x^2 - 1}$ e $\eu^{-t} = x - \sqrt{x^2 - 1}$, então
\[
T_n(x) = \cosh(nt)
= \frac{(x + \sqrt{x^2-1})^n + (x - \sqrt{x^2-1})^n}2 .
\]
Para $ x > 1$ o primeiro termo excede $\frac12(1)^n $ estritamente e
cresce geometricamente: $ T_n(x) > 1$.

\textbf{7.} De Moivre: $\cos n\theta = \Re\bigl((\cos\theta +
\iu\sin\theta)^n\bigr) = \sum_{2j \leq n}\binom n{2j}
\cos^{n-2j}\theta\,(\iu\sin\theta)^{2j}$, e $(\iu\sin\theta)^{2j}
= (-\sin^2\theta)^j = (\cos^2\theta - 1)^j $. Substituindo $ x =
\cos\theta $ e invocando a questão 3:
\[
T_n(x) = \sum_{0\leq 2j\leq n}\binom n{2j}x^{n-2j}(x^2 - 1)^j .
\]
Para $ n = 3$: $\binom30 x^3 + \binom32 x(x^2 - 1) = x^3 + 3x^3 -
3x = 4x^3 - 3x $, como na questão 1.

\textbf{8.} $ T_n(1) = \cos(n\cdot0) = 1$; $ T_n(-1) = \cos(n\pi) =
(-1)^n $; $ T_n(0) = \cos\frac{n\pi}2$, que é $0$ por estranho $ n $ e
$(-1)^{n/2}$ até mesmo $ n $.

\textbf{9.} Indução para $ U_n(\cos\theta) =
\frac{\sin(n+1)\theta}{\sin\theta}$: verdadeiro para $ U_0 = 1$ e $ U_1 =
2X $ ($\sin2\theta = 2\sin\theta\cos\theta $); o passo é o
identidade soma-produto $\sin(n+2)\theta = 2\cos\theta\,
\sin(n+1)\theta - \sin n\theta $. Agora diferencie
$ T_n(\cos\theta) = \cos n\theta $ em $\theta $:
$-\sin\theta\,T_n'(\cos\theta) = -n\sin n\theta $, então para $\theta
\notin \pi\Z $:
\[
T_n'(\cos\theta) = n\,\frac{\sin n\theta}{\sin\theta}
= n\,U_{n-1}(\cos\theta) ,
\]
e o \omterm{def:b1:poly:def}{polinômios} $ T_n'$ e $ nU_{n-1}$, concordando em
$\intoo{-1}1$, são iguais.

\textbf{10.} $\abs{\sin(n+1)\theta} = \abs{\sin n\theta\cos\theta
+ \cos n\theta\sin\theta} \leq \abs{\sin n\theta} +
\abs{\sin\theta}$, e a indução dá $\abs{\sin n\theta} \leq
n\abs{\sin\theta}$. Por isso $\abs{U_{n-1}} \leq n $ sobre $\intoo{-1}1$
e $\abs{T_n'} = n\abs{U_{n-1}} \leq n^2$ lá; no $\pm1$ o
limite se estende tomando limites (ou diretamente: $ U_{n-1}(1) = n $ de
a recorrência, $ U_n(1) = n + 1$ por indução, e a paridade dá
$ U_{n-1}(-1) = (-1)^{n-1}n $). Por isso $ T_n'(1) = n^2$ e $ T_n'(-1)
= (-1)^{n-1}n^2$: o limite $ n^2$ é alcançado nos extremos.

\textbf{11.} Diferenciar $\sin\theta\,T_n'(\cos\theta) = n\sin
n\theta $ (questão 9) em relação a $\theta $:
\[
\cos\theta\,T_n'(\cos\theta) - \sin^2\theta\,T_n''(\cos\theta)
= n^2\cos n\theta = n^2\,T_n(\cos\theta) .
\]
Com $ x = \cos\theta $ e $\sin^2\theta = 1 - x^2$: $ x\,T_n' - (1
- x^2)T_n'' = n^2T_n $ sobre $\intcc{-1}1$, portanto, em todos os lugares:
$(1 - x^2)y'' - xy' + n^2y = 0$ para $ y = T_n $. Verifique se há $ T_2 =
2x^2 - 1$: $(1 - x^2)(4) - x(4x) + 4(2x^2 - 1) = 4 - 4x^2 - 4x^2
+ 8x^2 - 4 = 0$.

\textbf{12.} $\widetilde T_n $ é \omterm{def:b1:poly:def}{Mônica} (questão 2) e
$\abs{\widetilde T_n} = 2^{1-n}\abs{T_n} \leq 2^{1-n}$ sobre
$\intcc{-1}1$, com $\widetilde T_n(y_k) = (-1)^k2^{1-n}$ no
$ n + 1$ pontos $ y_k $ (\cref{exo:b1:poly:10}): a norma é exatamente
$2^{1-n}$, obtido com sinais alternados.

\textbf{13.} $\widetilde T_n $ e $ P $ são ambos \omterm{def:b1:poly:def}{Mônica} de grau
$ n $, então os termos iniciais cancelam: $\deg D \leq n - 1$. No $ y_k $:
$ D(y_k) = (-1)^k2^{1-n} - P(y_k)$, e $\abs{P(y_k)} \leq \norm
P_\infty < 2^{1-n}$ força o sinal de $ D(y_k)$ ser aquele de
$(-1)^k2^{1-n}$, estritamente.

\textbf{14.} $ D $ muda de sinal entre $ y_{k+1}$ e $ y_k $ para cada
$ k = 0, \dots, n-1$: pela propriedade de valor intermediário, $ D $ tem um
raiz em cada um desses $ n $ intervalos abertos disjuntos entre pares --- $ n $
raízes distintas para um diferente de zero polinomial de grau $\leq n - 1$,
impossível. E $ D = 0$ também é impossível (as normas diferem).
Contradição: não \omterm{def:b1:poly:def}{Mônica} $ P $ de grau $ n $ tem $\norm P_\infty <
2^{1-n}$, que é o teorema de Chebyshev.

\textbf{15.} Agora $\abs{P(y_k)} \leq 2^{1-n}$ só, então $(-1)^k
D(y_k) = 2^{1-n} - (-1)^kP(y_k) \geq 2^{1-n} - \abs{P(y_k)} \geq
0$. Suponha $ D(y_k) = 0$ em um interior $ y_k $ ($0 < k < n $): então
$ P(y_k) = (-1)^k2^{1-n}$, então $\abs P $ atinge o seu supremo
$2^{1-n}$ no ponto interior $ y_k $, de onde $ P'(y_k) = 0$
(extremo interior); e $ T_n'(y_k) = nU_{n-1}(y_k) = 0$ desde
$\sin(n\cdot\frac{k\pi}n) = 0$ --- então $\widetilde T_n'(y_k) = 0$
também, e $ D'(y_k) = 0$: $ y_k $ é uma raiz de $ D $ de \omterm{def:b1:poly:derivative}{multiplicidade} pelo menos $2$.

\textbf{16.} Contar raízes de $ D $ com \omterm{def:b1:poly:derivative}{multiplicidade}. Deixar $ z $ seja o
número de pontos interiores $ y_k $ com $ D(y_k) = 0$ (cada um duplo
raiz, pela questão 15) e $ e \in \{0, 1, 2\}$ o número de
pontos finais ($ y_0$ ou $ y_n $) com $ D = 0$ (cada um uma raiz simples em
menos). Uma lacuna $(y_{k+1}, y_k)$ cujos dois pontos finais têm $ D
\neq 0$ carrega sinais estritamente alternados, daí um interior
raiz. Cada ponto interior que desaparece estraga no máximo seus dois
lacunas adjacentes, cada ponto final de fuga no máximo uma lacuna: pelo menos
$ n - 2z - e $ lacunas ainda contribuem com uma raiz cada, todas distintas
do $ y $-raízes. Total: pelo menos $(n - 2z - e) + 2z + e = n $
raízes com \omterm{def:b1:poly:derivative}{multiplicidade}, por um polinomial de grau $\leq n - 1$:
então $ D = 0$ e $ P = \widetilde T_n $. O minimizador é único.

\textbf{17.} O afim \omterm{def:b1:logic:map}{mapa} $ t \mapsto x = \frac{a+b}2 +
\frac{b-a}2\,t $ é uma bijeção $\intcc{-1}1 \to \intcc ab $. Se $ P $
é \omterm{def:b1:poly:def}{Mônica} de grau $ n $, então $ Q(t) = P(x(t))$ é um polinomial em
$ t $ com coeficiente líder $\bigl(\frac{b-a}2\bigr)^n $, e
$\sup_{\intcc ab}\abs P = \sup_{\intcc{-1}1}\abs Q $. O \omterm{def:b1:poly:def}{Mônica}
\omterm{def:b1:poly:def}{polinômio} $ Q/\bigl(\frac{b-a}2\bigr)^n $ tem sup-norma $\geq
2^{1-n}$ (questões 13-14), então
\[
\sup_{\intcc ab}\abs P \geq \Bigl(\frac{b-a}2\Bigr)^n 2^{1-n}
= 2\Bigl(\frac{b-a}4\Bigr)^n ,
\]
com igualdade exatamente para $ P(x) = \bigl(\frac{b-a}2\bigr)^n
\widetilde T_n\bigl(t(x)\bigr)$ (questão 16).

\textbf{18.} $\widetilde T_3 = \frac{T_3}4 = X^3 - \frac34X $;
$\widetilde T_3{}' = 3X^2 - \frac34$ desaparece em $\pm\frac12$.
Valores: $\widetilde T_3(-1) = -\frac14$, $\widetilde
T_3(-\tfrac12) = \frac14$, $\widetilde T_3(\tfrac12) = -\frac14$,
$\widetilde T_3(1) = \frac14$: quatro extremos alternados de
valor absoluto $\frac14$ --- então $\norm{\widetilde T_3}_\infty =
\frac14$, e pelo teorema de Chebyshev não \omterm{def:b1:poly:def}{Mônica} cúbico tem menor
sup-norma ativada $\intcc{-1}1$.

\textbf{19.} $\omega $ é \omterm{def:b1:poly:def}{Mônica} de grau $ n + 1$, então
$\norm\omega_\infty \geq 2^{-n}$ pelo teorema de Chebyshev (grau
$ n+1$), com igualdade se e somente se $\omega = \widetilde T_{n+1}
= 2^{-n}T_{n+1}$ (questão 16), ou seja, se e somente se os nós
são os $ n + 1$ raízes de $ T_{n+1}$. Com nós Chebyshev o erro
fator $\norm\omega_\infty $ é $2^{-n}$ --- o mínimo possível.

\textbf{20.} $5 \times 36^\circ = 180^\circ $, então $ T_5(c) =
\cos180^\circ = -1$: $16c^5 - 20c^3 + 5c + 1 = 0$. Teste $ x =
-1$: $-16 + 20 - 5 + 1 = 0$, e a expansão confirma
\[
16x^5 - 20x^3 + 5x + 1 = (x + 1)\bigl(4x^2 - 2x - 1\bigr)^2 .
\]
Desde $ c = \cos36^\circ \neq -1$, $ c $ é uma raiz de $4x^2 - 2x -
1$, cujas raízes são $\frac{1 \pm \sqrt5}4$; como $ c > 0$,
\[
\cos36^\circ = \frac{1 + \sqrt5}4 .
\]
Consistência: $\cos72^\circ = T_2(c) = 2c^2 - 1 = 2\cdot\frac{3 +
\sqrt5}8 - 1 = \frac{\sqrt5 - 1}4$, o valor encontrado em
\cref{exo:b1:complex:8}.

\textbf{21.} $\sqrt{1.1^2 - 1} = \sqrt{0.21} \approx 0.458$, então
$ x + \sqrt{x^2-1} \approx 1.558$ e $(1.558)^{10} \approx 84.5$,
enquanto $(1.1 - 0.458)^{10} \approx 0.01$: $ T_{10}(1.1) \approx
\frac{84.5 + 0.01}2 \approx 42$. UM polinomial preso em
$\intcc{-1}1$ no intervalo já passou $40$ um
décimo além de sua borda: o limite em um segmento não diz nada e
polegada fora dele.

\textbf{22.} Na fórmula da questão 7 para $ T_p $, O termo $ j = 0$
é $ X^p $; todos os outros termos carregam $\binom p{2j}$ com $0 < 2j <
p $ (observação $2j \neq p $ desde $ p $ é ímpar), que é divisível por $ p $
pelo primeiro passo da prova de \cref{thm:b1:arith:fermat}.
Portanto, todo coeficiente de $ T_p - X^p $ é um múltiplo de $ p $.
Verificações: $ T_3 - X^3 = 3X^3 - 3X = 3(X^3 - X)$; $ T_5 - X^5 = 15X^5
- 20X^3 + 5X = 5(3X^5 - 4X^3 + X)$.

\textbf{23.} Pela questão 17 com $\intcc ab = \intcc01$ e $ n =
2$: desvio mínimo $2\bigl(\frac14\bigr)^2 = \frac18$,
alcançado por $\bigl(\frac12\bigr)^2\widetilde T_2(2x - 1) =
\frac14\bigl((2x-1)^2 - \frac12\bigr) = x^2 - x + \frac18$. O
\omterm{def:b1:poly:def}{Mônica} quadrático mais próximo de zero em $\intcc01$ é $ x^2 - x +
\frac18$, com sup-norma $\frac18$.

\textbf{24.} (i) Rigidez --- a polinomial com mais raízes do que
seu grau é zero --- alimenta o princípio da singularidade
(questão 3), a transferência de identidades trigonométricas para
polinomial identidades (questões 4, 7, 9, 11), e ambas
argumentos de contagem de raízes da prova de extremalidade (questões 14,
16). (ii) A trigonometria de \cref{ch:b1:complex} (de Moivre,
soma ao produto) e o funções hiperbólicas de
\cref{ch:b1:functions} forneceu todas as identidades por trás da família;
a substituição $ x = \cos\theta $ é a ponte. (iii) O
\omterm{def:b1:arith:divides}{divisibilidade} $ p \mid \binom p{2j}$ de \cref{ch:b1:arith} virou
a fórmula do coeficiente no \omterm{def:b1:arith:congruence}{congruência} da questão 22.

\textbf{25.} O teorema de Chebyshev converte uma otimização sobre um
família de dimensão infinita (todos \omterm{def:b1:poly:def}{polinômios} mônicos) em finito
combinatória: um concorrente \omterm{def:b1:arith:prime}{melhor} que $\widetilde T_n $ faria
diferem dele por um baixo grau polinomial forçado a mudar de sinal
$ n $ vezes --- uma raiz a mais do que seu grau permite. O
padrão de equioscilação não é, portanto, uma curiosidade, mas a própria
certificado de otimalidade, e o caso de igualdade fica mais nítido
contagem de raízes com \omterm{def:b1:poly:derivative}{multiplicidades}. A substituição $ x =
\cos\theta $ merece a última palavra: transporta o rígido,
mundo discreto de \omterm{def:b1:poly:def}{polinômios} no mundo periódico de
trigonometria, onde raízes e extremos de $ T_n $ são simplesmente os
grade regular de $\cos n\theta$. Os dois fatos de análise emprestados
--- a propriedade de valor intermediário (questão 14; comprovada em
\cref{ch:b1:continuity}) e a derivada nula em um
extremo interior (questão 15; comprovado em \cref{ch:b1:derivative})
--- são exatamente as ferramentas que os capítulos posteriores devolverão,
fechando o ciclo.
\end{solution}
